\scanpage{11}

\uline{Géométries d'incidences et ensembles à opérateurs}

Soit \inserted{$\Pi=$} $(S,A,F\,;R,R')$ \deleted{\unsure{$P$}} une géométrie d'incidences de dim 2, \deleted{On appelle drapeau} \inserted{$\unsure{\Phi|P|\text{cf.}}$ l'ens ordonné} qu'il définit -- de ``dim. combinatoire'' 2 -- On appelle \uline{drapeau} \inserted{de $\Pi$} une partie tot. ordonnée non vide de $\Pi$ -- si son cardinal est $l$, on dit que le drapeau est de \uline{longueur} $l-1$. Le maximum des longueurs des drapeaux est 2 -- les drapeaux de longueur 2 sont appelés \uline{repères}, leur ens se note $\operatorname{Rep}(\Pi)$

\[
 \operatorname{Rep}(\Pi)\unsure{:|=}\{(\underset{\deleted{\unsure{0}}}{s},a,f)\in S\times A\times F\mid s<a<f\}
\]

Il y a des drapeaux des types suivants

\diagram{$\left.\begin{array}{ccc}0&1&2\\S&A&F\end{array}\right\}
 \begin{array}{l}\text{drapeaux minimaux \inserted{(de long. 0)}}\\\text{ou ``facettes''}\end{array}$}

\diagram{$\left.\begin{array}{ccc}
 (0,1)&(1,2)&(0,2)\\
 R\subset S\times A)&R'\subset A\times F&R''\subset S\times F
 \end{array}\right\}\ \text{drapeaux de long. 1}$}

\diagram{$\left.\begin{array}{c}(0,1,2)\\\operatorname{R\acute ep}\subset S\times A\times F\end{array}\right\}\ 
 \begin{array}{l}\text{repères, ou drapeaux}\\\text{maximaux}\end{array}$}

Il y a des applications naturelles ``d'oubli'' (partiel) entre ces ens. de drapeaux, en particulier

% archive sheet number: 10

